\section{Conclusion}

For similarity edges, Shannon redundancy with the node features follows analytically from
the construction, not from any experiment. What we tested, rather than assumed, is the
practical question -- whether this similarity-based audio graph nevertheless helps a real,
finite model -- across four tasks spanning text-only prediction, graph-only classification,
multimodal fusion, and cross-modal retrieval. Tasks 2 and 4 include a
direct, test-set graph-vs-edge-free comparison; Task 3's own such comparison is
validation-only, so its evidence instead comes from two test-set diagnostics: branch-zeroing
inside the trained fusion model, and a cross-tabulation of independently-trained
single-branch checkpoints. No \emph{learned-aggregation} graph-based model ever
showed a clear advantage over its edge-free alternative, and it trailed in Task 2's
pooling-matched reading -- a single-run difference, with GTZAN's small unresolved leakage
caveat and no uncertainty estimate, so we call this a trailing result, not a statistically
established loss.

One qualifier matters here. Task 2's non-learned exact 2-hop control, whose aggregation is
fixed rather than trained, does clearly beat every edge-free number -- a genuine complication,
discussed in Section~\ref{sec:discussion}, not swept into "graph-based model" for
convenience. Task 3's own standalone comparison was mixed and inconclusive on its own; in this
evidence base, branch-zeroing and complementarity -- an inference-time intervention on the
trained fusion model, and a correlational comparison of two separately-trained single-branch
models, respectively -- both point to little graph reliance instead. Section~\ref{sec:discussion}
also reports that the same two diagnostics, run against a second implementation's own
checkpoints, found substantially more reliance -- an unresolved disagreement between
implementations, not a second confirmation of the first. Task 4 showed no consistent advantage either way across three
seeds; with only three runs and no formal equivalence test, that is an absence of detected
difference, not proof of none.

We report the resulting conclusion as the paper's central finding: no task in this study
showed a clearly established graph advantage, though the supporting evidence and its strength
differ by task and are stated precisely rather than smoothed into one uniform result.
Presence of a graph branch in a music-understanding pipeline is not the same as demonstrated
reliance on it -- the two should be checked separately, not assumed to go together.
